\chapter{Experimentación}

A lo largo de este capítulo se recopilará la experimentación realizada con \textit{JaTyC}. Recordemos que el propósito principal es analizar el poder y las limitaciones de la herramienta, y a su vez validar su comportamiento correcto en casos interesantes. También buscamos entender mejor el potencial de la verificación de protocolos de objetos con typestates en general.

Los tests fueron realizados con pequeños extractos de código, compuestos por diferentes clases, interfaces, protocolos y clientes. Estos se analizaron con la herramienta y se anotó el comportamiento observado. Cada sección del capítulo tratará un aspecto diferente de \textit{JaTyC} o verificación de protocolos de objetos en general, e incluirá algunas de las pruebas más relevantes.

\section{Protocolos Compuestos}

Es de interés saber si \textit{JaTyC} tiene en consideración el estado de los colaboradores internos de un objeto protocolado. La herramienta parecería interpretar un typestate símplemente como un conjunto de métodos habilitados para una instancia y sus transiciones hacia otros typestates. El estado interno de un objeto pasa a ser secundario a la herramienta, aunque entiende que solo debe de ser modificado por métodos protocolados y no los \textit{anytime}. En este sentido, el protocolo y uso debido de un objeto está dictado solamente por las secuencias de sus métodos válidas.

Pero podría el typestate de un objeto ser dictado también por el estado de un colaborador interno? Es decir, podría el protocolo de un objeto padre ser uno compuesto donde, dependiendo del typestate de un colaborador, llamadas a un mismo método resulten en distintos resultados (tal como excepciones en ocasiones)? Buscamos entender cómo la herramienta analiza estos casos, donde la complejidad de determinar el comportamiento correcto de un objeto sobresale del scope del mismo.

\subsection{Incumplimiento de Protocolo}

Nuestras primeras pruebas realizadas buscaban romper el protocolo de un colaborador interno, mientras se respetaba el del objeto padre. Como mencionamos, era de interés saber si, según el punto de vista de un cliente, se podían provocar comportamientos inesperados en un objeto a pesar de utilizarlo correctamente, causados por su estado interno.

\begin{figure}[htbp]
    \centering
    \lstinputlisting[language=Java]{./code/capitulo4/collaborator_compound_typestate/Robot.java}
    \caption{Extracto de código de la clase \texttt{Robot}.}
    \label{exp:compound-1:robot}
\end{figure}

\begin{figure}[htbp]
    \centering
    \lstinputlisting{./code/capitulo4/collaborator_compound_typestate/RobotProtocol.protocol}
    \caption{Protocolo de la clase \texttt{Robot}.}
    \label{exp:compound-1:robot-p}
\end{figure}

El código y protocolo de una clase \texttt{Robot} puede observarse en las Figuras \ref{exp:compound-1:robot} y \ref{exp:compound-1:robot-p}. El robot se mueve en un mapa de tamaño $2\times2$, una unidad de distancia a la vez con métodos \texttt{moveLeft()}, \texttt{moveRight()}, \texttt{moveUp()} y \texttt{moveDown()}. Comienza en la posición inicial arriba a la izquierda y de moverse a una posición afuera del mapa levantará una excepción, simulando ``caerse del mapa''. Su protocolo refleja este comportamiento, con un typestate para cada posición posible del robot y solo movimientos permitidos tal que nunca se caiga.

Un primer experimento con este robot fue realizado con el controlador de las Figuras \ref{exp:compound-1:controller} y \ref{exp:compound-1:controller-p}. El \texttt{RobotController} cuenta con un \texttt{Robot} de colaborador interno, y simplemente toma instrucciones de un cliente para moverlo. Su protocolo es pequeño, contando con un método de control inicial \texttt{canMoveRobot()}, el cual habilita al método \texttt{move(String direction)}, que mueve el robot en la dirección deseada.

\begin{figure}[htbp]
    \centering
    \lstinputlisting[language=Java]{./code/capitulo4/collaborator_compound_typestate/RobotController.java}
    \caption{Código de la clase \texttt{RobotController}.}
    \label{exp:compound-1:controller}
\end{figure}

\begin{figure}[htbp]
    \centering
    \lstinputlisting{./code/capitulo4/collaborator_compound_typestate/RobotControllerProtocol.protocol}
    \caption{Protocolo de la clase \texttt{RobotController}.}
    \label{exp:compound-1:controller-p}
\end{figure}

Finalmente, el código cliente de la Figura \ref{exp:compound-1:error} sigue el protocolo del controlador pero rompe el del robot. Esto es porque, estando en la posición inicial, lo mueve hacia la izquierda donde cae fuera del mapa.

Observando el error devuelto por la herramienta, claramente detecta que se incumple el protocolo del robot. De todos modos, no se reporta al código cliente como culpable, sino al método \texttt{move(string direction)} del \texttt{RobotController}. Esto tiene sentido, dado que puede mover al robot por fuera del mapa, por ejemplo hacia arriba o a la izquierda estando en la posición inicial. El controlador no utiliza métodos de control antes de mover el robot para asegurarse de que es seguro hacerlo.

\begin{figure}[htbp]
    \centering
    \lstinputlisting[language=Java]{./code/capitulo4/collaborator_compound_typestate/Main.java}
    \vspace{0.3cm}
    \lstinputlisting{./code/capitulo4/collaborator_compound_typestate/Error.txt}
    \caption{Errores reportados por \textit{JaTyC} al analizar el código del controlador en la Figura \ref{exp:compound-1:controller} y su respectivo código cliente \texttt{Main}.}
    \label{exp:compound-1:error}
\end{figure}

Una conclusión que puede sacarse de esto es que el verdadero cliente del robot no es el código en \texttt{Main}, sino el controlador mismo. \textit{JaTyC} reportó el posible incumplimiento del protocolo desde el \texttt{RobotController} porque es allí donde el robot se utiliza. Por ende, no analiza ambos protocolos como uno compuesto desde \texttt{Main}, sino que verifica a cada cliente de un objeto protocolado por separado.

Una variación del test que demuestra el mismo comportamiento se puede ver en las Figuras \ref{exp:compound-3:controller} y \ref{exp:compound-3:controller-p}. Mientras que el \texttt{Robot} no tiene cambios, el \texttt{RobotController} ahora cuenta con un método \texttt{run()}, que hace el robot de una vuelta al mapa en el sentido del reloj, y \texttt{moveRight()} que mueve el robot hacia la derecha una posición. Notemos que la única posición segura del robot desde la cual puede llamarse a \texttt{run()} es la inicial.

\begin{figure}[htbp]
    \centering
    \lstinputlisting[language=Java]{./code/capitulo4/collaborator_compound_typestate_3/RobotController.java}
    \caption{Código de la clase \texttt{RobotController}.}
    \label{exp:compound-3:controller}
\end{figure}

\begin{figure}[htbp]
    \centering
    \lstinputlisting{./code/capitulo4/collaborator_compound_typestate_3/RobotControllerProtocol.protocol}
    \caption{Protocolo de la clase \texttt{RobotController}.}
    \label{exp:compound-3:controller-p}
\end{figure}

\textit{JaTyC} analizó el código con dos clientes diferentes. El primero solamente realizaba una llamada a \texttt{canMoveRobot()} y \texttt{run()}, repetando ambos protocolos del robot y controlador. El segundo (Figura \ref{exp:compound-3:error}) llama a \texttt{run()}, luego a \texttt{moveRight()} y finalmente a \texttt{run()}, todas precedidas por llamdas a \texttt{canMoveRobot()}. Comenzando en la posición inicial del mapa $2x2$ arriba a la izquierda, la primera llamada a \texttt{run()} deja al robot en la misma posición. Al llamar a \texttt{moveRight()} el robot pasa a estar en la posición arriba derecha. Finalmente, la última llamada a \texttt{run()} provoca que el robot se caiga del mapa, dado que comienza con un movimiento a la derecha, y rompa su protocolo.

\begin{figure}[htbp]
    \centering
    \lstinputlisting[language=Java]{./code/capitulo4/collaborator_compound_typestate_3/Main.java}
    \vspace{0.3cm}
    \lstinputlisting{./code/capitulo4/collaborator_compound_typestate_3/Error.txt}
    \caption{Errores reportados por \textit{JaTyC} al analizar el código del controlador en la Figura \ref{exp:compound-3:controller} y su respectivo código cliente \texttt{Main}.}
    \label{exp:compound-3:error}
\end{figure}

Para ambos códigos clientes, la herramienta reporta los mismos errores de la Figura \ref{exp:compound-3:error}. Los mismos reportan sólamente sobre el uso indebido del robot por parte del controlador, donde diferentes combinaciones de llamadas de métodos provocan que caiga del mapa. Esto nuevamente confirma que para \textit{JaTyC} los protocolos no son compuestos, y se analiza por separado el uso del controlador por el código cliente, y del robot por parte del controlador. 

% Esto es respaldado por lo visto en la sección TODO, donde el algoritmo de análisis de clases repasa todas las posibles combinaciones de llamadas de métodos para una clase particular, realizando un seguimiento de los posibles estados para parámentros, valores de retorno y sus colaboradores internos.

\subsection{Finalización de Protocolo}

Por otro lado, realizamos experimentaciones enfocadas en la finalización de protocolos compuestos. Es decir, cómo maneja la herramienta que uno de los dos protocolos termine antes que el otro. En el caso del objeto padre terminando su protocolo primero, sabrá \texttt{JaTyC} que el colaborador interno no podrá finalizar el suyo? Y de ser al revéz, sabrá que el objeto padre no puede utilizar los métodos protocolados del colaborador?

El primero de los casos puede observarse en la Figura . Las clases son muy similares a las mencionados anteriormente, con algunas diferencias. El \texttt{Robot} ahora tiene un método \textit{anytime}, \texttt{dance()}, y la única forma de finalizar su protocolo es a través del nuevo método \texttt{rest()}. El \texttt{RobotController} cuenta con un nuevo método \texttt{turnOff()}, que representa la única forma de transicionar a \texttt{end}, y desde ningún método llama a \texttt{rest()} del robot (en otras palabras, no sabe finalizar su protocolo).

En la Figura se puede ver el código cliente analizado. En el mismo, se instancia un conto
% TODO me falta corregir un poco estos tests.
\section{Correctitud Vs Orden de Métodos}

El análisis de los object protocols en general se limita a verificar el orden correcto de las llamadas a métodos. Mientras que \textit{JaTyC} puede hacer mucho más que eso, como detectar posibles errores de puntero nulo o especificar typestates esperados para parámetos o valores de retorno, no tiene conocimiento sobre el contenido de los méotodos por sobre los que predica. En otras palabras, no juzga la \textbf{correctitud} de los métodos de un objeto protocolado, y por ende pueden suceder excepciones inesperadas incluso aunque se siga su contrato. En ese sentido, es interesante indagar cómo la herramienta está limitada en este aspecto, y cómo podemos utilizarla para hacer que nuestro código se alinee con lo su especificación de correcto.

En la Figura puede verse uno de los tests que realizamos, con una clase \texttt{Stack} desarrollada incorrectamente. Observando su protocolo, se puede apilar y desapilar elementos mientras no esté llena ni vacía, respectivamente. Por ejemplo, solamente se puede llamar a \texttt{push()} si \texttt{isFull()} retorna falso. Para esta implementación, el método \texttt{isFull()} siempre devuelve falso, con lo cual permite que se apile hasta sobrepasar el tamaño del stack y levantar una excepción. El código cliente usa la pila siguiendo su protocolo, pero apilando más elementos que su capacidad.

Al correr \textit{JaTyC}, efectivamente no reporta ningún error, pero al ejecutar el código cliente termina con una excepción. Esto muestra que, incluso respetando el protocolo de un objeto y utilizando sus métodos de control (como \texttt{isFull()}) antes de sus métodos críticos (como \texttt{push()}), \textit{JaTyC} no nos asegura que esté implementado de forma correcta.

Por otro lado, de cambiar el método \texttt{isFull()} para que correctamente controle cuándo se puede apilar, el \texttt{Stack} y su protocolo es un ejemplo de cómo podemos apoyarnos en la herramienta para facilitar que el código sea más correcto.


%   - faulty_stack 
%     - Un stack que su implementación tiene fallas y larga excepciones, pero su protocolo lo cumple bien.
%   - normal_stack
%     - Esto es un stack normal, donde para hacer que cumpla un protocolo hay que agregarle métodos booleanos de control. Sirve para explicar que la tool no puede trackear cantidad de llamadas a métodos o secuencias de métodos, a no ser que hagamos un estado por llamada. Es decir, en escencia lo que se testea es querer determinar cuándo estoy por ejemplo haciendo un push cuando la pila ya está llena. En ese sentido, la herramienta al no tener una pila interna es incapaz de asegurarse que el código sea seguro, y te obliga a usar métodos booleanos de control.
%     - Explicar que si no puede probar que tu código es seguro, te fuerza a hacer checkeos (por medio de métodos de control) para que lo sea. En ese sentido, son falsos positivos, donde potencialmente el codigo esta bien pero te obliga a ser mas defensivo.

\section{Intervención Humana}

\section{Subtipado}

\section{Estado Vs Typestate}

\section{Casos Bordes} % o \section{Capacidad de Modelado}

\section{Aliasing}

\section{Excepciones}